% GENERATED FILE. Edit canonical lessons, then run npm run generate:books.
% canonical-source-hash: fnv1a64:1866489cb035a769
% canonical-lessons: JA-C419-kokuseki, JA-C419-nenrei, JA-C419-seibetsu, JA-C419-shimei, JA-C419-uketsukeru, JA-R419-review-414-419

\chapter{Name, Age and Nationality}
\label{ch:ja-name-age-and-nationality}
\hlchaptermodality{419}

\begin{chapteropening}
\textbf{By the end of this chapter:} \emph{I can give my details at a counter.}

Complete the last lesson of chapter 419: I can give my details at a counter.
\end{chapteropening}

\section[\texorpdfstring{\ja{こくせき} (kokuseki)}{こくせき (kokuseki)}]{\ja{こくせき} (kokuseki) — nationality}

\label{lesson:JA-C419-kokuseki}



\begin{quote}
Before the new one: say the Japanese for a box on a form, then the Japanese for a date.
\end{quote}

\subsection*{You'll want to know: \ja{こくせき}}
\textbf{\ja{こくせき}} — \emph{kokuseki} — \textquotedblleft{}nationality\textquotedblright{}.

\begin{quote}
\textbf{\ja{ようしに} \ja{こくせきを} \ja{きにゅうしました。}} — \emph{yōshi ni kokuseki o kinyū shimashita} — \textquotedblleft{}I filled in my nationality on the form.\textquotedblright{}
\end{quote}

\begin{cousinweb}
\textbf{\ja{こくせき}} (\emph{kokuseki}) is Sino-Japanese (kango): two Chinese characters read with their on'yomi, readings borrowed from Chinese. They are \emph{koku}, \textquotedblleft{}country\textquotedblright{} (Middle Chinese \emph{kwok}) and \emph{seki}, \textquotedblleft{}domiciliary register\textquotedblright{} (\emph{dzjek}) (KANJIDIC2; Middle Chinese in Baxter's transcription).
\end{cousinweb}

\begin{grammarlens}[title={What you've built}]
One more word for getting around.
\end{grammarlens}

\subsection*{Your turn}
\begin{itemize}
\raggedright
  \item \emph{Say it:} \emph{kokuseki}
  \item \emph{Say it:} \emph{kokuseki}, once more
  \item \emph{Say it:} \emph{yōshi ni kokuseki o kinyū shimashita} — I filled in my nationality on the form.
  \item \emph{Recall:} say the Japanese for a box on a form, then the Japanese for a date, then say \emph{kokuseki} again
\end{itemize}

\begin{tcolorbox}[breakable,colback=teal!4,colframe=teal!35!black,title={Before you move on}]
What does \ja{こくせき} mean? (\textquotedblleft{}Nationality\textquotedblright{}.) Say it once more.
\end{tcolorbox}
\section[\texorpdfstring{\ja{ねんれい} (nenrei)}{ねんれい (nenrei)}]{\ja{ねんれい} (nenrei) — age (a formal word)}

\label{lesson:JA-C419-nenrei}



\begin{quote}
Before the new one: say the Japanese for a date, then the Japanese for nationality.
\end{quote}

\subsection*{You'll want to know: \ja{ねんれい}}
\textbf{\ja{ねんれい}} — \emph{nenrei} — \textquotedblleft{}age (a formal word)\textquotedblright{}.

\begin{quote}
\textbf{\ja{ここに} \ja{ねんれいを} \ja{かいて} \ja{ください。}} — \emph{koko ni nenrei o kaite kudasai} — \textquotedblleft{}Please write your age here.\textquotedblright{}
\end{quote}

\begin{cousinweb}
\textbf{\ja{ねんれい}} (\emph{nenrei}) is Sino-Japanese (kango): two Chinese characters read with their on'yomi, readings borrowed from Chinese. They are \emph{nen}, \textquotedblleft{}year\textquotedblright{} (Middle Chinese \emph{nen}) and \emph{rei}, \textquotedblleft{}age\textquotedblright{} (\emph{leng}) (KANJIDIC2; Middle Chinese in Baxter's transcription).
\end{cousinweb}

\begin{grammarlens}[title={What you've built}]
One more word for getting around.
\end{grammarlens}

\subsection*{Your turn}
\begin{itemize}
\raggedright
  \item \emph{Say it:} \emph{nenrei}
  \item \emph{Say it:} \emph{nenrei}, once more
  \item \emph{Say it:} \emph{koko ni nenrei o kaite kudasai} — Please write your age here.
  \item \emph{Recall:} say the Japanese for a date, then the Japanese for nationality, then say \emph{nenrei} again
\end{itemize}

\begin{tcolorbox}[breakable,colback=teal!4,colframe=teal!35!black,title={Before you move on}]
What does \ja{ねんれい} mean? (\textquotedblleft{}Age (a formal word)\textquotedblright{}.) Say it once more.
\end{tcolorbox}
\section[\texorpdfstring{\ja{せいべつ} (seibetsu)}{せいべつ (seibetsu)}]{\ja{せいべつ} (seibetsu) — sex (on a form)}

\label{lesson:JA-C419-seibetsu}



\begin{quote}
Before the new one: say the Japanese for nationality, then the Japanese for age (a formal word).
\end{quote}

\subsection*{You'll want to know: \ja{せいべつ}}
\textbf{\ja{せいべつ}} — \emph{seibetsu} — \textquotedblleft{}sex (on a form)\textquotedblright{}.

\begin{quote}
\textbf{\ja{せいべつの} \ja{らんも} \ja{あります。}} — \emph{seibetsu no ran mo arimasu} — \textquotedblleft{}There is a box for sex as well.\textquotedblright{}
\end{quote}

\begin{cousinweb}
\textbf{\ja{せいべつ}} (\emph{seibetsu}) is Sino-Japanese (kango): two Chinese characters read with their on'yomi, readings borrowed from Chinese. They are \emph{sei}, \textquotedblleft{}sex\textquotedblright{} (Middle Chinese \emph{sjengH}) and \emph{betsu}, \textquotedblleft{}separate\textquotedblright{} (\emph{bjet}) (KANJIDIC2; Middle Chinese in Baxter's transcription).
\end{cousinweb}

\begin{grammarlens}[title={What you've built}]
One more word for getting around.
\end{grammarlens}

\subsection*{Your turn}
\begin{itemize}
\raggedright
  \item \emph{Say it:} \emph{seibetsu}
  \item \emph{Say it:} \emph{seibetsu}, once more
  \item \emph{Say it:} \emph{seibetsu no ran mo arimasu} — There is a box for sex as well.
  \item \emph{Recall:} say the Japanese for nationality, then the Japanese for age (a formal word), then say \emph{seibetsu} again
\end{itemize}

\begin{tcolorbox}[breakable,colback=teal!4,colframe=teal!35!black,title={Before you move on}]
What does \ja{せいべつ} mean? (\textquotedblleft{}Sex (on a form)\textquotedblright{}.) Say it once more.
\end{tcolorbox}
\section[\texorpdfstring{\ja{しめい} (shimei)}{しめい (shimei)}]{\ja{しめい} (shimei) — one's full name}

\label{lesson:JA-C419-shimei}



\begin{quote}
Before the new one: say the Japanese for age (a formal word), then the Japanese for sex (on a form).
\end{quote}

\subsection*{You'll want to know: \ja{しめい}}
\textbf{\ja{しめい}} — \emph{shimei} — \textquotedblleft{}one's full name\textquotedblright{}.

\begin{quote}
\textbf{\ja{しめいと} \ja{じゅうしょを} \ja{かいて} \ja{ください。}} — \emph{shimei to jūsho o kaite kudasai} — \textquotedblleft{}Please write your full name and address.\textquotedblright{}
\end{quote}

\begin{cousinweb}
\textbf{\ja{しめい}} (\emph{shimei}) is Sino-Japanese (kango): two Chinese characters read with their on'yomi, readings borrowed from Chinese. They are \emph{shi}, \textquotedblleft{}family name\textquotedblright{} (Middle Chinese \emph{dzyeX}) and \emph{mei}, \textquotedblleft{}name\textquotedblright{} (\emph{mjieng}) (KANJIDIC2; Middle Chinese in Baxter's transcription).
\end{cousinweb}

\begin{grammarlens}[title={What you've built}]
One more word for getting around.
\end{grammarlens}

\subsection*{Your turn}
\begin{itemize}
\raggedright
  \item \emph{Say it:} \emph{shimei}
  \item \emph{Say it:} \emph{shimei}, once more
  \item \emph{Say it:} \emph{shimei to jūsho o kaite kudasai} — Please write your full name and address.
  \item \emph{Recall:} say the Japanese for age (a formal word), then the Japanese for sex (on a form), then say \emph{shimei} again
\end{itemize}

\begin{tcolorbox}[breakable,colback=teal!4,colframe=teal!35!black,title={Before you move on}]
What does \ja{しめい} mean? (\textquotedblleft{}One's full name\textquotedblright{}.) Say it once more.
\end{tcolorbox}
\section[\texorpdfstring{\ja{うけつける} (uketsukeru)}{うけつける (uketsukeru)}]{\ja{うけつける} (uketsukeru) — to accept applications; to take bookings}

\label{lesson:JA-C419-uketsukeru}



\begin{quote}
Before the new one: say the Japanese for sex (on a form), then the Japanese for one's full name.
\end{quote}

\subsection*{You'll want to know: \ja{うけつける}}
\textbf{\ja{うけつける}} — \emph{uketsukeru} — \textquotedblleft{}to accept applications; to take bookings\textquotedblright{}.

\begin{quote}
\textbf{\ja{まどぐちは} \ja{ごじまで} \ja{よやくを} \ja{うけつけます。}} — \emph{madoguchi wa goji made yoyaku o uketsukemasu} — \textquotedblleft{}The counter takes bookings until five.\textquotedblright{}
\end{quote}

Polite \textbf{\ja{うけつけます}}, \emph{uketsukemasu}; an ichidan verb. Past (polite): \textbf{\ja{まどぐちが} \ja{ようしを} \ja{うけつけました。}}, \emph{madoguchi ga yōshi o uketsukemashita}.

\begin{cousinweb}
\textbf{\ja{うけつける}} (\emph{uketsukeru}) is a native Japanese word (wago), inherited rather than borrowed from Chinese, built from \emph{uke}, \textquotedblleft{}receiving\textquotedblright{} (from \emph{ukeru}) and \emph{tsukeru}, \textquotedblleft{}to attach\textquotedblright{}; its kanji were attached later as kun'yomi (JMdict).
\end{cousinweb}

\begin{grammarlens}[title={What you've built}]
One more word for getting around.
\end{grammarlens}

\subsection*{Your turn}
\begin{itemize}
\raggedright
  \item \emph{Say it:} \emph{uketsukeru}
  \item \emph{Say it:} \emph{uketsukeru}, once more
  \item \emph{Say it:} \emph{madoguchi wa goji made yoyaku o uketsukemasu} — The counter takes bookings until five.
  \item \emph{Recall:} say the Japanese for sex (on a form), then the Japanese for one's full name, then say \emph{uketsukeru} again
\end{itemize}

\begin{tcolorbox}[breakable,colback=teal!4,colframe=teal!35!black,title={Before you move on}]
What does \ja{うけつける} mean? (\textquotedblleft{}To accept applications; to take bookings\textquotedblright{}.) Say it once more.
\end{tcolorbox}
\section[(dialogue)]{Looking back: from Hot, Damp Days to Name, Age and Nationality}

\label{lesson:JA-R419-review-414-419}



\begin{quote}
Say the Japanese for humidity and damp and humid.
\end{quote}

\subsection*{Your turn}
Say these aloud:

\begin{itemize}
\raggedright
  \item humidity, damp and humid, the sun's rays, shade, dazzling
  \item a sudden evening shower, light rain, soaking wet, to get something wet, sheltering from the rain
  \item dawn, sunrise, a sunset glow, to turn very cold, to shine on
  \item hand luggage, a luggage tag, to carry on one's back, to carry in one's arms, to wheel along
  \item a signature, a blank form, filling in a form, a box on a form, a date
  \item nationality, age (a formal word), sex (on a form), one's full name, to accept applications
\end{itemize}

\begin{tcolorbox}[breakable,colback=teal!4,colframe=teal!35!black,title={Before you move on}]
Humidity? (\textbf{\ja{しつど}}, \emph{shitsudo}.) Damp and humid? (\textbf{\ja{じめじめ}}, \emph{jimejime}.) The sun's rays? (\textbf{\ja{ひざし}}, \emph{hizashi}.) Shade? (\textbf{\ja{ひかげ}}, \emph{hikage}.) Dazzling? (\textbf{\ja{まぶしい}}, \emph{mabushii}.) A sudden evening shower? (\textbf{\ja{ゆうだち}}, \emph{yūdachi}.) Light rain? (\textbf{\ja{こさめ}}, \emph{kosame}.) Soaking wet? (\textbf{\ja{びっしょり}}, \emph{bisshori}.) To get something wet? (\textbf{\ja{ぬらす}}, \emph{nurasu}.) Sheltering from the rain? (\textbf{\ja{あまやどり}}, \emph{amayadori}.) Dawn? (\textbf{\ja{よあけ}}, \emph{yoake}.) Sunrise? (\textbf{\ja{ひので}}, \emph{hinode}.) A sunset glow? (\textbf{\ja{ゆうやけ}}, \emph{yūyake}.) To turn very cold? (\textbf{\ja{ひえこむ}}, \emph{hiekomu}.) To shine on? (\textbf{\ja{てらす}}, \emph{terasu}.) Hand luggage? (\textbf{\ja{てにもつ}}, \emph{tenimotsu}.) A luggage tag? (\textbf{\ja{にふだ}}, \emph{nifuda}.) To carry on one's back? (\textbf{\ja{せおう}}, \emph{seou}.) To carry in one's arms? (\textbf{\ja{かかえる}}, \emph{kakaeru}.) To wheel along? (\textbf{\ja{ころがす}}, \emph{korogasu}.) A signature? (\textbf{\ja{しょめい}}, \emph{shomei}.) A blank form? (\textbf{\ja{ようし}}, \emph{yōshi}.) Filling in a form? (\textbf{\ja{きにゅう}}, \emph{kinyū}.) A box on a form? (\textbf{\ja{らん}}, \emph{ran}.) A date? (\textbf{\ja{ひにち}}, \emph{hinichi}.) Nationality? (\textbf{\ja{こくせき}}, \emph{kokuseki}.) Age (a formal word)? (\textbf{\ja{ねんれい}}, \emph{nenrei}.) Sex (on a form)? (\textbf{\ja{せいべつ}}, \emph{seibetsu}.) One's full name? (\textbf{\ja{しめい}}, \emph{shimei}.) To accept applications? (\textbf{\ja{うけつける}}, \emph{uketsukeru}.)
\end{tcolorbox}
